\begin{example}[General Linear group]
    Let us consider \(G = \mathrm{GL}(n, \mathbb{R}) = \{M \in \mathbb{R}^{n\times n} : \det(M) \neq 0\}\), we are interested in finding the Lie algebra of the grop as the tangent space at the identity element \(e = \mathbb{I}_n\) of the group. We can identify the elements of the group with points in \(\mathbb{R}^{n^2}\) by mapping each matrix \(M\) to its entries \(M^{AB}\). Locally, around the identity element, the tangent space can be identified with the space of all \(n \times n\) matrices:
    \[
        T_e G \simeq \mathfrak{g} = T_e \mathbb{R}^{n \times n} \simeq \mathbb{R}^{n^2},
    \]
    and since the dimension of the algebra, the group, the manifold etc. coincides, we know it to be of dimension \(n^2\). When passing to the complex counterpart, we have \(G = \mathrm{GL}(n, \mathbb{C})\) and the Lie algebra is \(\mathfrak{g} = T_e G \simeq \mathbb{C}^{n^2}\), which is of dimension \(2n^2\) as a real vector space.
\end{example}

\begin{example}[Special Linear group]
    If we consider the special linear group \(G = \mathrm{SL}(n, \mathbb{R}) = \{M \in \mathrm{GL}(n, \mathbb{R}) : \det(M) = 1\}\), we can identify the elements of the group with points in \(\mathbb{R}^{n^2}\) by mapping each matrix \(M\) to its entries \(M^{AB}\). Let \(M(t) : \mathbb{R} \to G\) be a curve in the group with \(M(0) = e = \mathbb{I}_n\). Then the tangent vector at the identity is given by
    \[
        v = \frac{\mathrm{d}}{\mathrm{d} t} M(t) \bigg|_{t=0}.
    \]
    Since \(M(t)\) is a matrix in \(\mathrm{SL}(n, \mathbb{R})\), we have \(\det(M(t)) = 1\) for all \(t\). Taking the derivative of this condition with respect to \(t\) and evaluating at \(t = 0\) gives
    \[
        \frac{\mathrm{d}}{\mathrm{d} t} \det(M(t)) \bigg|_{t=0} = 0.
    \]
    Using the formula for the derivative of the determinant
    \[
        \frac{\mathrm{d}}{\mathrm{d} t} \det(M(t)) \bigg|_{t=0} = \Tr\left[\det(M(t)) (M^{-1}(t))^T \dot{M}(t)\right]_{t=0} = \Tr(\dot{M}(0)) = \Tr(v),
    \]
    this implies
    \[
        \mathrm{tr}(v) = 0,
    \]
    which means that the tangent vector \(v\) (derivative of an element) at the identity element must be a traceless matrix. Therefore, the Lie algebra of the special linear group is the set of all traceless \(n \times n\) matrices, which can be identified with a subspace of \(\mathbb{R}^{n^2}\):
    \[
        \mathfrak{sl}(n, \mathbb{R}) = T_e G \simeq \{X \in \mathbb{R}^{n \times n} : \mathrm{tr}(X) = 0\}.
    \]
    The dimension of this Lie algebra is \(n^2 - 1\), since the condition of being traceless reduces the number of independent parameters by one (but is the same as the original group, since there we had the condition on the determinant). In the complex case, we have \(G = \mathrm{SL}(n, \mathbb{C})\) and the Lie algebra is \(\mathfrak{sl}(n, \mathbb{C}) = T_e G \simeq \{X \in \mathbb{C}^{n \times n} : \mathrm{tr}(X) = 0\}\), which is of dimension \(2(n^2 - 1)\) as a real vector space.
\end{example}

\begin{example}[Orthogonal group]
    Consider the orthogonal group \(G = \mathrm{O}(n, \mathbb{R}) = \{M \in \mathbb{R}^{n \times n} : M^T M = \mathbb{I}_n\}\). To find its Lie algebra, we look at the tangent space at the identity element \(e = \mathbb{I}_n\). A curve \(M(t) : \mathbb{R} \to G\) with \(M(0) = \mathbb{I}_n\) satisfies \(M(t)^T M(t) = \mathbb{I}_n\) for all \(t\). Differentiating this condition with respect to \(t\) and evaluating at \(t = 0\) gives
    \[
        \frac{\mathrm{d}}{\mathrm{d} t} (M(t)^T M(t)) \bigg|_{t=0} = 0.
    \]
    Using the product rule,
    \[
        \dot{M}(0)^T \mathbb{I}_n + \mathbb{I}_n^T \dot{M}(0) = 0,
    \]
    which simplifies to
    \[
        \dot{M}(0)^T + \dot{M}(0) = 0.
    \]
    This means that the tangent vector \(v = \dot{M}(0)\) must be skew-symmetric. Therefore, the Lie algebra of the orthogonal group is the set of all skew-symmetric \(n \times n\) matrices:
    \[
        \mathfrak{o}(n, \mathbb{R}) = T_e G \simeq \{X \in \mathbb{R}^{n \times n} : X^T + X = 0\}.
    \]
    The dimension of this Lie algebra is \(n(n-1)/2\), as there are \(n(n-1)/2\) independent parameters for a skew-symmetric matrix.
\end{example}

\begin{example}[Special Orthogonal group]
    For the special orthogonal group \(G = \mathrm{SO}(n, \mathbb{R}) = \{M \in \mathrm{O}(n, \mathbb{R}) : \det(M) = 1\}\), the Lie algebra is the same as that of the orthogonal group, since the condition of having determinant equal to one does not impose any additional constraints on the tangent vectors at the identity element. Therefore, we have
    \[
        \mathfrak{so}(n, \mathbb{R}) = T_e G \simeq \{X \in \mathbb{R}^{n \times n} : X^T + X = 0\},
    \]
    which is of dimension \(n(n-1)/2\).

    We can study this group as the intersection of the special linear group and the orthogonal group, since it is defined by the conditions of being both orthogonal and having determinant equal to one. Thus we can write
    \[
        \mathrm{SO}(n, \mathbb{R}) = \mathrm{SL}(n, \mathbb{R}) \cap \mathrm{O}(n, \mathbb{R}),
    \]
    and consequently we can write the Lie algebra as the intersection of the Lie algebras of the special linear group and the orthogonal group:
    \[
        \mathfrak{so}(n, \mathbb{R}) = \mathfrak{sl}(n, \mathbb{R}) \cap \mathfrak{o}(n, \mathbb{R}) = \{X \in \mathbb{R}^{n \times n} : X^T + X = 0, \mathrm{tr}(X) = 0\}.
    \]
    But since the condition of being skew-symmetric already implies that the trace is zero, we can simplify this to
    \[
        \mathfrak{so}(n, \mathbb{R}) = \{X \in \mathbb{R}^{n \times n} : X^T + X = 0\} = \mathfrak{o}(n, \mathbb{R}),
    \]
    which is the same as the Lie algebra of the orthogonal group.
\end{example}

\begin{example}[Unitary group]
    For the unitary group \(G = \mathrm{U}(n, \mathbb{C}) = \{U \in \mathbb{C}^{n \times n} : U^\dagger U = \mathbb{I}_n\}\), we can find its Lie algebra by looking at the tangent space at the identity element \(e = \mathbb{I}_n\). A curve \(U(t) : \mathbb{R} \to G\) with \(U(0) = \mathbb{I}_n\) satisfies \(U(t)^\dagger U(t) = \mathbb{I}_n\) for all \(t\). Differentiating this condition with respect to \(t\) and evaluating at \(t = 0\) gives
    \[
        \frac{\mathrm{d}}{\mathrm{d} t} (U(t)^\dagger U(t)) \bigg|_{t=0} = 0.
    \]
    Using the product rule,
    \[
        \dot{U}(0)^\dagger \mathbb{I}_n + \mathbb{I}_n^\dagger \dot{U}(0) = 0,
    \]
    which simplifies to
    \[
        \dot{U}(0)^\dagger + \dot{U}(0) = 0.
    \]
    This means that the tangent vector \(v = \dot{U}(0)\) must be anti-Hermitian. Therefore, the Lie algebra of the unitary group is the set of all anti-Hermitian \(n \times n\) matrices:
    \[
        \mathfrak{u}(n, \mathbb{C}) = T_e G \simeq \{X \in \mathbb{C}^{n \times n} : X^\dagger + X = 0\}.
    \]
    The real dimension of this Lie algebra is \(n^2\), as there are \(n^2\) independent parameters for an anti-Hermitian matrix: the upper half of the matrix has \(\frac{n(n-1)}{2}\) independent parameters, and the diagonal elements (all purely imaginary since they have to be changed in sign when daggered) have \(n\) independent parameters, for a total of \(2 \frac{n(n-1)}{2} + n = n^2\) independent parameters.
\end{example}

\begin{example}[Special Unitary group]
    For the special unitary group \(G = \mathrm{SU}(n, \mathbb{C}) = \{U \in \mathrm{U}(n, \mathbb{C}) : \det(U) = 1\}\), the Lie algebra is the set of all anti-Hermitian \(n \times n\) matrices with trace zero:
    \[
        \mathfrak{su}(n, \mathbb{C}) = T_e G \simeq \{X \in \mathbb{C}^{n \times n} : X^\dagger + X = 0, \mathrm{tr}(X) = 0\}.
    \]
    The real dimension of this Lie algebra is \(n^2 - 1\), as the condition of being traceless reduces the number of independent parameters by one.
\end{example}

We are computing the \textbf{real dimension} of the lie algebra even in the complex case, since we are considering the Lie algebra as a real vector space. If we were to consider it as a complex vector space, we would have to consider that when the coefficients are complex, the dimension would be halved, since each complex parameter corresponds to two real parameters (the real and imaginary parts). Thus, when we have an odd number of independent parameters, the complex dimension would be a half-integer, which is not possible for a vector space. Therefore, we consider the Lie algebra as a real vector space, even in the complex case, to ensure that the dimension is always an integer.

We want to impose only \textbf{holomorphic functions} on the complex Lie groups, since we want to preserve the complex structure of the group and its Lie algebra. If we were to allow non-holomorphic functions, we would lose the complex structure and we would not be able to define a complex Lie algebra. The holomorphic functions (for example the determinant consition, which is just a polynomial equation in complex variables) ensure that the Lie algebra is closed under the Lie bracket and that it has a well-defined complex structure, which is essential for many applications in mathematics and physics.

\begin{lemma}
    Let \(G\) be a matrix Lie group, and let \(v \in T_h G\) be a tangent vector in \(h\) and consider \(a \in G\); then the pushforward of \(v\) under the left translation by \(a\) is given by
    \[
        [[(L_a)_*]_h(v)]^{AB} = \sum_{C=1}^n a^{AC} v^{CB} = (a \cdot v)^{AB},
    \]
    where the last equality makes sense only in a matrix group, since we are using the matrix multiplication to define the action of the group on the tangent vector. This shows how the pushforward of a tangent vector under left translation can be expressed in terms of the matrix multiplication of the group element and the tangent vector, which is a key property that allows us to understand how the Lie algebra is related to the structure of the Lie group.
\end{lemma}
\begin{proof}
    We can use a test function \(\phi \in \mathbb{C}^\infty(G)\) to compute the pushforward of \(v\) under left translation by \(a\):
    \[
        [(L_a)_*]_h(v)(\phi) = v(\phi \circ L_a) = \sum_{k} v^k \partial_k (\phi \circ L_a) \big|_{x(h)} \equiv v(\tilde{\phi}),
    \]
    where we defined \(\tilde{\phi} := \phi (ag)\). Now we can express the function \(\tilde{\phi}\) in terms of the matrix entries \(g^{AB}\) if we compute
    \[
        [[(L_a)_*]_h(v)]^{AB} = v^{AB} \frac{\partial \phi}{\partial g^{PQ}} \bigg|_{g = ah} \frac{\partial}{\partial g^{AB}} (ag)^{PQ} \bigg|_{g = h},
    \]
    which can be rewritten as
    \[
        = v^{AB} \frac{\partial \phi}{\partial g^{PQ}} \bigg|_{g = ah} a^{PR} \delta^{RA}\delta^{QB} = a^{PR} v^{RQ} \frac{\partial \phi}{\partial g^{PQ}} = (a \cdot v)^{AB} \frac{\partial \phi}{\partial g^{AB}},
    \]
    [...]
\end{proof}

\begin{proposition}
    Let \(G\) be the matrix group and \(v,w \in T_e G \implies \mathfrak{g}\), then we can associate to \(v\) and \(w\) the left-invariant vector fields \(V^{(v)}\) and \(W^{(w)}\), and we can compute the Lie bracket of these vector fields at the identity element:
    \[
        [V^{(v)}, W^{(w)}]_e = [v,w]_\mathfrak{g},
    \]
    where the Lie bracket on the right-hand side is defined by the commutator of the corresponding tangent matrices:
    \[
        [v,w]_\mathfrak{g} = (v w)^{AB} - (w v)^{AB} = v^{AC} w^{CB} - w^{AC} v^{CB},
    \]
    which is the standard Lie bracket for matrix Lie algebras and makes sense in the context of matrix groups, since the tangent vectors can be identified with matrices and the Lie bracket can be defined as the commutator of these matrices.
\end{proposition}

We are not going to prove this proposition, since it is a straightforward computation that follows from the definition of the Lie bracket of vector fields and the properties of the pushforward under left translations. The key point is that the left-invariant vector fields are defined by translating the tangent vectors at the identity element to every other point in the group using left translations, and thus their Lie bracket can be computed by using the properties of the pushforward under left translations and the definition of the Lie bracket of vector fields.

\section{Exponential Map}

We are now going to define the exponential map for a Lie group, which is a crucial tool for understanding the relationship between the Lie algebra and the Lie group. The exponential map allows us to relate elements of the Lie algebra (tangent vectors at the identity element) to elements of the Lie group, and it provides a way to construct curves in the Lie group from tangent vectors in the Lie algebra. It is the tool which let us take a finite transformation from an infinitesimal one, and thus it is fundamental for applications in physics, where we often deal with infinitesimal transformations and we want to understand how they relate to finite transformations.

\begin{definition}[Integral curve]
    Let \(V \in \Gamma(T \mathcal{M})\) be a vector field and let \(\gamma: (a,b) \to \mathcal{M}\) be an integral curve for \(V\), i.e.,a smooth curve whose tangent vector at each point \(\gamma(t)\) is equal to the value of the vector field at that point \(V_{\gamma(t)}\), i.e.
    \[
        \dot{\gamma}(t) = V_{\gamma(t)} \quad \forall t \in (a,b).
    \]
    If we use a chart \((U,x)\) around a point \(p \in \mathcal{M}\) such that \(\gamma(t_0) = p\) for some \(t_0 \in (a,b)\), then we can write the integral curve in coordinates as
    \[
        \dot{\gamma}^i(t) = \frac{\mathrm{d}}{\mathrm{d}t} \gamma^i(t) = V^i(\gamma(t)) \quad \forall t \in (a,b), \quad \text{(ODE)},
    \]
    which is a system of ordinary differential equations (ODEs) that can be solved to find the integral curve \(\gamma(t)\) given the vector field \(V\) and the initial condition \(\gamma(t_0) = p\).
\end{definition}

The existence and uniqueness of solutions to this ODE is guaranteed by the standard theory of ODEs, and thus we can always find an integral curve for any given vector field and initial condition.

\begin{definition}[One-parameter subgroup]
    Let a smooth curve \(\sigma : \mathbb{R} \to G\) in a Lie group \(G\) is called a \textbf{one-parameter subgroup} of \(G\) if it satisfies the group homomorphism property
    \[
        \sigma(s+t) = \sigma(s) \sigma(t) \quad \forall s,t \in \mathbb{R},
    \]
    with \(\sigma(0) = e\) is the identity element of the Lie group \(G\).
\end{definition}

This last definition means that a one-parameter subgroup is a curve in the Lie group that respects the group structure, i.e., the value of the curve at the sum of two parameters is equal to the product (group operation) of the values of the curve at each parameter. This property is crucial for understanding how one-parameter subgroups relate to the Lie algebra and how they can be used to construct finite transformations from infinitesimal ones.

We will adopt the notation in which
\[
    \sigma \equiv \{\sigma(t) \mid t \in \mathbb{R}\},
\]
as a subgroup of \(G\) generated by the one-parameter subgroup \(\sigma\). This notation emphasizes the fact that the one-parameter subgroup generates a subgroup of \(G\) that is isomorphic to \(\mathbb{R}\) under addition, and it allows us to think of the one-parameter subgroup as a continuous family of group elements parameterized by a real variable \(t\).

\begin{proposition}
    Let \(G\) be a Lie group and let \(V \in \Gamma(TG)\) be a left-invariant vector field on \(G\). Then the integral curve \(\sigma : \mathbb{R} \to G\) of \(V\) with initial condition \(\sigma(0) = e\) is a one-parameter subgroup of \(G\).
\end{proposition}

The idea here is that by taking the integral curve of a left-invariant vector field starting at the identity element, we obtain a curve that respects the group structure and thus defines a one-parameter subgroup. The left-invariance of the vector field ensures that the flow generated by the vector field is compatible with the group operation, which is essential for the resulting curve to be a one-parameter subgroup.

\begin{proposition}
    Let \(\sigma\) be a one-parameter subgroup of a Lie group \(G\). Then the tangent vector \(V_{\sigma(t)} = \frac{d}{dt}\sigma(t)\) to the curve \(\sigma(t)\) at any point \(t\) can be computed as the pushforward of the tangent vector at the identity element \(V_e\) under the left translation by \(\sigma(t)\):
    \[
        V_{\sigma(t)}  = [(L_{\sigma(t)})_*]_e (\dot{\sigma}(0)),
    \]
    which implies that \(V_{\sigma(t)}\) is a left-invariant vector field along the curve \(\sigma(t)\), and thus the one-parameter subgroup \(\sigma\) is generated by the left-invariant vector field \(V\) that corresponds to the tangent vector at the identity element.
\end{proposition}

Thus every one-parameter subgroup of a Lie group is generated by a left-invariant vector field, and conversely, every left-invariant vector field generates a one-parameter subgroup. This establishes a deep connection between the Lie algebra (the space of left-invariant vector fields) and the structure of the Lie group (the one-parameter subgroups), which is fundamental for understanding how the Lie algebra encodes information about the local structure of the Lie group and how it can be used to construct finite transformations from infinitesimal ones.

It is a one-to-one correspondence between the elements of the Lie algebra and the one-parameter subgroups of the Lie group, which is a key result in the theory of Lie groups and Lie algebras.

\begin{corollary}
    Suppose to have two different one-parameter subgroups \(\sigma(t)\) and \(\mu(t)\) of a Lie group \(G\) with \(\dot{\sigma}(0) = \dot{\mu}(0) = v \in T_e G\). Then \(\sigma(t) = \mu(t)\) for all \(t \in \mathbb{R}\).
\end{corollary}

This corollary can be read as saying that the one-parameter subgroup generated by a given tangent vector at the identity element is unique. This emphasizes the one-to-one correspondence between the elements of the Lie algebra and the one-parameter subgroups of the Lie group, and it ensures that each element of the Lie algebra generates a unique one-parameter subgroup.

\begin{definition}
    Let \(X \in \mathfrak{g}\) be an element of the Lie algebra of a Lie group \(G\). We denote as \(\sigma_X\) the unique one-parameter subgroup of \(G\) constructed as follows:
    \begin{enumerate}
        \item \(X \in \mathfrak{g} \iff V^{(X)} \in \Gamma(TG)\) is left-invariant vector field on \(G\) such that \(V^{(X)}_e = X\).
        \item \(\sigma_X : \mathbb{R} \to G\) is the unique integral curve of \(V^{(X)}\) with initial condition \(\sigma_X(0) = e\) (passing through the identity element) with \(\dot{\sigma}_X(0) = V_e^{(X)} = X\); it is a one-parameter subgroup of \(G\).
    \end{enumerate}
\end{definition}

Since along the curve \(\sigma_X(t)\) we have a possible mapping from the Lie algebra to the Lie group, we can define the exponential map as follows.

\begin{definition}[Exponential map]
    The exponential map
    \[
        \exp : \mathfrak{g} \to G, \quad X \mapsto \sigma_X(1),
    \]
    is called the exponential map of the Lie group \(G\). It maps an element \(X\) of the Lie algebra \(\mathfrak{g}\) to the value of the one-parameter subgroup \(\sigma_X\) at \(t = 1\), which is an element of the Lie group \(G\).
\end{definition}

\begin{proposition}
    If \(X \in \mathfrak{g}\), \(\lambda \in \mathbb{R}\) then we have
    \[
        \exp(\lambda X) = \sigma_X(\lambda).
    \]
    Since this is a one-parameter subgroup, we have
    \[
        \exp((\lambda + \mu) X) = \sigma_X(\lambda + \mu) = \sigma_X(\lambda) \sigma_X(\mu) = \exp(\lambda X) \exp(\mu X), \quad \forall \lambda, \mu \in \mathbb{R},
    \]
    and we can also derive straightforwardly that \(\exp(0) = e\) and \(\exp(-X) = \sigma_X(-1) = \sigma_X(1)^{-1} = \exp(X)^{-1}\). This shows that the exponential map has properties that are consistent with the group structure, and it allows us to relate the Lie algebra to the Lie group in a way that respects the algebraic structure of both.
\end{proposition}

This definition allows us to plug in the map any vector from the Lie algebra and obtain an element of the Lie group, which is a powerful tool for understanding how the Lie algebra encodes information about the local structure of the Lie group and how it can be used to construct finite transformations from infinitesimal ones.

\begin{proof}
    For \(\lambda \neq 0\) we can define a new curve \(\tilde{\sigma} : \mathbb{R} \to G\) as
    \[
        \tilde{\sigma}(t) := \sigma_X(\lambda t),
    \]
    which is still a one-parameter subgroup of \(G\) since it satisfies the group homomorphism property. If we now define another one-parameter subgroup from \(\sigma_{\lambda X}(0) = \lambda X \), then we have also
    \[
        \dot{\tilde{\sigma}}(0) = \frac{\mathrm{d}}{\mathrm{d} t} \sigma_X(\lambda t) \bigg|_{t=0} = \lambda \dot{\sigma}_X(0) = \lambda X,
    \]
    which means that \(\tilde{\sigma}\) and \(\sigma_{\lambda X}\) are two one-parameter subgroups of \(G\) with the same tangent vector at the identity element (initial conditions), and thus by the previous corollary they must coincide, i.e., \(\tilde{\sigma}(t) = \sigma_{\lambda X}(t)\) for all \(t \in \mathbb{R}\). If this is the case, then we can use the definition of the exponential map with respect to the one-parameter subgroup \(\sigma_{\lambda X}\) to write
    \[
        \exp(\lambda X) = \sigma_{\lambda X}(1) = \tilde{\sigma}(1) = \sigma_X(\lambda),
    \]
    which is the desired result.
\end{proof}

We are now going to express some basic and important properties of the exponential map.

\begin{proposition}
    The exponential map has the following properties:
    \begin{enumerate}
        \item The image of the exponential map is a neighborhood of the identity element \(e\) in \(G\), which means that
              \[
                  \exp(\mathfrak{g}) \subseteq G_0 = \{g \in G | \exists \gamma : [a,b] \to G \text{ continuous, with } \gamma(t_0) = e, \gamma(t_1) = g\},
              \]
              where \(G_0\) is the connected component of the identity element in \(G\). For example \(\mathrm{SO}(n) = [O(n)]_e\).
        \item The Lie algebra \(\mathfrak{g}\) is a \(\mathbb{R}\)-vector space \(\simeq \mathbb{R}^n\), then we can say that the exponential map is a smooth map from \(\mathbb{R}^n\) to \(G\), and thus it is a \textbf{local diffeomorphism} around the identity element \(e\). This means that there exists a neighborhood \(U\) of \(0 \in \mathfrak{g}\) such that \(\exp : U \to \exp(U)\) is a diffeomorphism onto its image, which is a neighborhood of the identity element in \(G\). With these restrictions the exponential map is thus injective.
        \item The exponential map \(\exp : \mathfrak{g} \to G\) is a smooth map, but generally not surjective, since there may be elements of the Lie group that cannot be reached by exponentiating elements of the Lie algebra. But if \(G_0\) is a compact (connected component) then it is surjective. For example, if \(G = \mathrm{SL}(2, \mathbb{R})\), which is not compact, then the exponential map is not surjective. On the other hand, if \(G = \mathrm{SU}(n)\), which is compact, then the exponential map is surjective. There may be cases of surjectivity even for non-compact groups, but this is not guaranteed in general (for example, the exponential map for \(\mathrm{GL}(n,\mathbb{C})\) and \(\mathrm{SO}(1,n)\) is surjective even if the group is not compact).
        \item Any \(g \in G_0\) can be written as
              \[
                  g = \exp(X_1) \exp(X_2) \cdots \exp(X_k),
              \]
              for some \(X_1, X_2, \ldots, X_k \in \mathfrak{g}\) and some \(k \in \mathbb{N}\). This means that any element of the connected component of the identity can be expressed as a finite product of exponentials of elements of the Lie algebra, which do not contradicts the previous point about surjectivity, since we are not claiming that every element of \(G\) can be written as a single exponential, but rather that it can be written as a finite product of exponentials.
        \item \textbf{Baker-Campbell-Hausdorff (BCH) formula}: For any \(X,Y \in \mathfrak{g}\), we have then
              \[
                  Z = X + Y + \frac{1}{2} [X,Y]_\mathfrak{g} + \frac{1}{12} [X,[X,Y]_\mathfrak{g}]_\mathfrak{g} - \frac{1}{12} [Y,[X,Y]_\mathfrak{g}]_\mathfrak{g} + \cdots,
              \]
              where the series continues with higher-order commutators. This formula let us connect ...
              Furtheore, if this series converges, we have
              \[
                  \exp(Z) = \exp(X) \exp(Y) = \exp\left(X + Y + \frac{1}{2} [X,Y]_\mathfrak{g} + \frac{1}{12} [X,[X,Y]_\mathfrak{g}]_\mathfrak{g} - \frac{1}{12} [Y,[X,Y]_\mathfrak{g}]_\mathfrak{g} + \cdots\right),
              \]
              has a solution for \(Z\) in terms of \(X\) and \(Y\). This formula is fundamental for understanding how the group operation in \(G\) can be expressed in terms of the Lie algebra \(\mathfrak{g}\), and it provides a way to compute the product of two exponentials in terms of the Lie bracket and higher-order commutators. If the group is compact or the map is surjective, then we can use the BCH formula to express any element of the group as an exponential of an element of the Lie algebra (since the series would always converge in this case). For the case in which \([X,Y]=0\) then the BCH formula simplifies to
              \[
                  \exp(X) \exp(Y) = \exp(X + Y),
              \]
              which is a special case that occurs when the Lie algebra is abelian (i.e., when all elements commute with each other).
        \item Let \(F : G_1 \to G_2\) be a smooth map between two Lie groups (Lie group homomorphism), then
              \[
                  \exp_2(F_*(X)) = F(\exp_1(X)), \quad \forall X \in \mathfrak{g}_1,
              \]
              where \(\exp_1\) and \(\exp_2\) are the exponential maps of \(G_1\) and \(G_2\) respectively, and \(F_*\) is the pushforward of \(F\). This property shows that the exponential map is compatible with Lie group homomorphisms, and it allows us to relate the exponential maps of different Lie groups through the pushforward of a homomorphism between them.
    \end{enumerate}
\end{proposition}

\subsection{Will see}

[...]

\begin{proposition}[Exponential on matrix group]
    Let $G$ be a matrix group, then \(\forall X \in \mathfrak{g}\) we have
    \[
        [\exp(X)]^{AB} = (e^X)^{AB},
    \]
    where \(X=X^{AB}\) is the tangent matrix and
    \[
        e^X = \sum_{n=0}^\infty \frac{X^n}{n!},
    \]
    and this is exactly the exponential map defined in the previous section.
\end{proposition}

\begin{proof}
    Let \(\sigma_{X}\) be the integral curve of the left-invariant vector field associated with \(X\), \(V^{(X)}\); we have from the definition of integral curve that
    \[
        \dot{\sigma}_{X}(t) = V^{(X)}_{\sigma_{X}(t)} = [(L_{\sigma_{X}(t)})_{*}]_e (X).
    \]
    Now we can use the exponential map to write
    \[
        (\frac{\mathrm{d} }{\mathrm{d} t} \exp(tX))^{AB} = [\dot{\sigma}_{X}(t)]^{AB} = [\sigma_{X}(t)]^{AC} X^{CB} = [\exp(tX)]^{AC} X^{CB}.
    \]
    Since the one parameter subgroup \(\sigma_{X}\) satisfies the composition
    \[
        \sigma_{X}(t+s) = \sigma_{X}(t) \sigma_{X}(s),
    \]
    then if we fix one of the parameters, say \(s\), we can treat \(\sigma_{X}(s)\) as a constant matrix and solve the above ODE to get
    \[
        \frac{\mathrm{d}}{\mathrm{d} t} [\sigma_{X}(t+s)]^{AB}\bigg|_{s} = [\sigma_{X}(s)]^{AC} X^{CB} \implies [\sigma_{X}(s)]^{AB} = \frac{\mathrm{d} }{\mathrm{d} \tau} \sigma(\tau)^{AB} \bigg|_{\tau=t+s} = \left( [(L_{\sigma_{X}(t+s)})_*]_e (X) \right)^{AB},
    \]
    which can be read as the tangent vector of \(\sigma_{X}\) at the point \(\sigma_{X}(t+s)\) which is given by the pushforward at the identity. We have then:
    \[
        \frac{\mathrm{d}}{\mathrm{d} t} [\sigma_{X}(t+s)]^{AB}\bigg|_{s} = \sigma_X(t+s)^{AC} X^{CB} = \sigma_X(t)^{AC} \sigma_X(s)^{CD} X^{DB},
    \]
    since we are still making computations which makes sense only when speaking about matrices. Then
    \[
        \frac{\mathrm{d}}{\mathrm{d} t} [\sigma_{X}(t+s)]^{AB}\bigg|_{s} = \dots
    \]
    \TODO{substitute t + s with t + t2 at t1}
    When we consider \(t_1 = 0\) the term \(\sigma_{X}(t_1)\) is just the identity matrix, so we have
    \[
        (\sigma_X(t_2) X)^{AB} = (X \sigma_X(t_2))^{AB}, \quad \forall t_2 \in \mathbb{R},
    \]
    thus we have that \(\sigma_X(t)\) commutes with \(X\) for all \(t\). Then the above ODE can be solved by the power series expansion of the exponential function, and we have
    \[
        \implies (\frac{\mathrm{d} }{\mathrm{d} t} \exp(tX))^{AB} = (X \exp(tX))^{AB} = (\exp(tX) X)^{AB}.
    \]
    Given this equation we know that the solution is given by the power series expansion of the exponential function, and we needed to show the commutation of \(\sigma_X(t)\) with \(X\) to be able to solve the ODE (since ...).
\end{proof}